\documentclass[11pt,a4paper]{article}
\usepackage[utf8]{inputenc}
\usepackage[portuguese]{babel}
\usepackage{amsmath,amsfonts,amssymb}
\usepackage{graphicx}
\usepackage{booktabs}
\usepackage{hyperref}
\usepackage{geometry}
\usepackage{float}
\usepackage{listings}
\usepackage{xcolor}

\geometry{margin=25mm}

\definecolor{codegreen}{rgb}{0,0.6,0}
\definecolor{codegray}{rgb}{0.5,0.5,0.5}
\definecolor{codeblue}{rgb}{0.1,0.2,0.6}
\definecolor{backcolour}{rgb}{0.96,0.96,0.96}

\lstdefinestyle{mystyle}{
    backgroundcolor=\color{backcolour},   
    commentstyle=\color{codegreen},
    keywordstyle=\color{codeblue},
    numberstyle=\tiny\color{codegray},
    stringstyle=\color{purple},
    basicstyle=\ttfamily\footnotesize,
    breakatwhitespace=false,         
    breaklines=true,                 
    captionpos=b,                    
    keepspaces=true,                 
    numbers=left,                    
    numbersep=5pt,                  
    showspaces=false,                
    showstringspaces=false,
    showtabs=false,                  
    tabsize=2
}
\lstset{style=mystyle}

\title{\textbf{Relatório Técnico: Previsão de Preços de Imóveis com Fusão de Dados e Redes MLP}}
\author{\textbf{Felipe Silva Loschi} \\ Engenharia de Software (GES - 601)}
\date{\today}

\begin{document}

\maketitle

\begin{abstract}
Este relatório mostra o desenvolvimento de um modelo de Machine Learning para prever preços de casas usando características físicas do imóvel e dados demográficos da região. Juntando as duas bases de dados através do código postal, criei uma arquitetura de rede neural Multi-Layer Perceptron (MLP) com TensorFlow e Keras. O pipeline foi feito com cuidado para evitar vazamento de dados (\textit{Data Leakage}) e trata problemas não-lineares aplicando transformação logarítmica e uma seleção de variáveis baseada no coeficiente de Spearman e em regras de negócio.
\end{abstract}

\section{Introdução e Fusão de Dados}
O objetivo do projeto é criar um modelo capaz de prever o preço de venda de imóveis residenciais. Para deixar o modelo mais inteligente, juntei a base com os dados físicos das casas (\texttt{kc\_house\_data.csv}) com o dataset de informações demográficas da região (\texttt{zipcode\_demographics.csv}), usando a coluna \texttt{zipcode} como chave para o \textit{merge}.

Colunas que serviam apenas como identificadores ou dados temporais que não ajudariam na predição (\texttt{id} e \texttt{date}) foram retiradas do dataset. Além disso, criei uma nova variável binária chamada \texttt{was\_renovated} a partir do ano de reforma (\texttt{yr\_renovated}) para indicar simplesmente se a casa foi reformada ou não, simplificando essa informação para a rede neural.

\section{Análise Exploratória de Dados (EDA) e Seleção de Atributos}
Na análise exploratória, estudei como as variáveis preditoras se distribuem e como elas se relacionam com o preço (\texttt{price}). Como muitos dados não seguem uma distribuição linear perfeita e possuem assimetria, usei o coeficiente de correlação de \textbf{Spearman} para guiar a escolha das melhores variáveis.

\subsection{Identificação de Outliers e Filtragem}
Durante a EDA, encontrei um erro claro na coluna de quantidade de quartos (\texttt{bedrooms}): um dos imóveis estava registrado com mais de 30 quartos, o que não fazia sentido para o tamanho da casa. Fiz um filtro para remover esse outlier e manter apenas os registros reais ($<30$ quartos), evitando que esse dado inconsistente atrapalhasse o treino do modelo.

\subsection{Critério de Seleção de Atributos}
Para não sobrecarregar a rede MLP com variáveis irrelevantes ou redundantes, defini um critério de corte: selecionei apenas as variáveis que tinham uma correlação em módulo igual ou maior que $0,3$ ($|\rho| \ge 0.3$) com a variável alvo \texttt{price}. Esse valor é um bom limite prático para garantir que estamos usando apenas variáveis que realmente trazem alguma informação útil para o modelo.

\begin{figure}[H]
    \centering
    \includegraphics[width=1\textwidth]{heatmap_spearman.png}
    \caption{Matriz de Correlação de Spearman filtrada para a variável dependente Price.}
    \label{fig:heatmap}
\end{figure}

Apesar de seguir esse corte matemático, abri duas exceções importantes pensando na lógica do negócio:
\begin{itemize}
    \item \textbf{\texttt{view}:} A variável que mede a qualidade da vista da casa teve correlação de $0,29$. Como ficou muito perto do corte de $0,3$ e a vista é um fator que pesa muito na decisão de quem vai comprar um imóvel, decidi manter essa coluna no modelo.
    \item \textbf{\texttt{long}:} Mantive a longitude para fazer par com a latitude (\texttt{lat}). Se eu tirasse a longitude, o modelo perderia a noção de geolocalização exata, ficando impossibilitado de entender quais bairros ou regiões específicas são mais valorizados.
\end{itemize}

Também retirei as colunas com a contagem absoluta de escolaridade (\texttt{edctn\_*\_qty}) deixando apenas as respectivas variáveis percentuais, que estão em uma escala comum a todos. No final, o modelo ficou com 17 variáveis preditoras.

\begin{figure}[H]
    \centering
    \includegraphics[width=1\textwidth]{grid_features_selecionadas.png}
    \caption{Grid de dispersão das 17 features selecionadas após o corte de correlação e ajustes de domínio.}
    \label{fig:grid_features}
\end{figure}

\subsection{Transformação Logarítmica}
Notei que o preço (\texttt{price}) e outras variáveis financeiras/demográficas, como renda (\texttt{medn\_incm\_per\_prsn\_amt}) e valor médio de moradia regional (\texttt{hous\_val\_amt}), tinham distribuições muito esticadas para a direita (presença de valores muito altos). Para resolver isso e ajudar o algoritmo a convergir melhor, apliquei a transformação logarítmica natural:
\begin{equation}
    \tilde{x} = \ln(x)
\end{equation}

Essa transformação ajuda a deixar os dados em uma escala mais equilibrada, facilitando o ajuste dos pesos da rede neural durante o treino.

\begin{figure}[H]
    \centering
    \includegraphics[width=1\textwidth]{boxplot_features.png}
    \caption{Análise de Dispersão e Limites de Outliers via Boxplots das Features Selecionadas.}
    \label{fig:boxplots}
\end{figure}

\section{Engenharia de Recursos e Pipeline de Preparação}
\subsection{Prevenção de Data Leakage}
Para garantir que o modelo funcione bem com dados novos e do mundo real (\texttt{future\_unseen\_examples.csv}), dividi a base em 80\% para treino e 20\% para teste. Na hora de padronizar os dados com o \texttt{StandardScaler}, fiz o processo do jeito correto para evitar o vazamento de dados (\textit{Data Leakage}): o ajuste (\textit{fit}) foi feito apenas no conjunto de treino, usando apenas o \textit{transform} nos dados de teste e validação:
\begin{lstlisting}[language=Python]
# Isolamento estrito do Scaler para evitar Data Leakage
standard = StandardScaler()
X_train = standard.fit_transform(X_train)
X_test = standard.transform(X_test)
\end{lstlisting}

\section{Arquitetura da Rede Neural MLP}
O modelo escolhido foi uma rede neural Multi-Layer Perceptron (MLP) criada com Keras e TensorFlow. Montei uma estrutura com camadas densas que vão diminuindo de tamanho para extrair os padrões dos dados de forma eficiente.

\begin{table}[H]
\centering
\caption{Especificação Arquitetural da Rede MLP Dedicada}
\label{tab:mlp}
\begin{tabular}{lcc}
\toprule
\textbf{Camada} & \textbf{Neurônios / Unidades} & \textbf{Função de Ativação} \\ \midrule
Input (Entrada) & 17 & --- \\
Dense 1         & 256 & ReLU \\
Dense 2         & 128 & ReLU \\
Dense 3         & 64 & ReLU \\
Dense 4         & 32 & ReLU \\
Saída (Output)  & 1 & Linear (Regressão) \\ \bottomrule
\end{tabular}
\end{table}

Usei a função de perda de Erro Quadrático Médio (\textit{MSE}) e o otimizador \textbf{Adam} com a taxa de aprendizado padrão ($\alpha = 10^{-3}$). Nas camadas escondidas usei a função de ativação ReLU para evitar problemas com o sumiço do gradiente, e na última camada usei ativação linear, que é o padrão para problemas de regressão.

\section{Análise de Resultados e Generalização}
Para evitar o problema de \textit{overfitting} e garantir a generalização do modelo, usei o \textit{train\_test\_split}, dividindo o conjunto em 80\% para treino e 20\% para teste. Além disso, usei o \textit{Early Stopping} para monitorar o erro de um conjunto de validação (20\% separado do treino), configurado com uma paciência de 15 épocas.

\begin{figure}[H]
      \centering
      \includegraphics[width=0.7\textwidth]{curva_convergencia.png}
      \caption{Histórico de Treinamento: Convergência das funções de perda no Treino e Validação.}
      \label{fig:loss}
\end{figure}

O comportamento das curvas de erro de treino e validação caindo juntas mostra que o modelo aprendeu os padrões de forma estável. No fim das contas, isso significa que o modelo está pronto para generalizar bem e pronto para ser testado na prática para estimar preços de imóveis de forma confiável.

\end{document}